\writeSectionFrame{\presentationTitle}{Das Kompositum-Muster}
\begin{frame}{Steckbrief}
	\begin{itemize}
 		\item Deutscher Name: Kompositum
 		\item Englischer Name: Composite
 		\item Gruppe: Strukturmuster
 	\end{itemize}
\end{frame}

\begin{frame}{Kompositum}
	\begin{block}{Zweck}
 		Füge Objekte zu Baumstrukturen zusammen, um Teil-Ganzes-Hierarchien zu repräsentieren. Das Kompositummuster ermöglicht es Klienten, einzelne Objekte sowie Kompositionen von Objekten einheitlich zu behandeln.
 	\end{block}
\end{frame}

\begin{frame}{Allgemeines Klassendiagramm}
	\begin{center}
		\begin{tikzpicture}[scale=0.7, transform shape]
		    \umlclass[x=0, y=0]{Komponente}{
		        -name: String
		    }{
		        +operation(): void
		    }
		
		    \umlclass[x=4, y=-4]{Blatt}{
		        -name: String
		    }{
		        +operation(): void
		    }
		
		    \umlclass[x=-4, y=-4]{Kompositum}{
		        -name: String \\
		        -elemente: List<Komponente>
		    }{
		        +operation(): void \\
		        +hinzufuegen(Komponente): void \\
		        +entfernen(Komponente): void \\
		        +getKind(int): Komponente
		    }
		
		    \umlinherit{Blatt}{Komponente}
		    \umlinherit{Kompositum}{Komponente}
		    \umlaggreg[geometry=|-, anchor1=north, anchor2=west]{Kompositum}{Komponente}
		\end{tikzpicture}
	 \end{center}
\end{frame}

\begin{frame}{Klassendiagramm - Beispiel Mailinglisten}
	\begin{center}
		\begin{tikzpicture}[scale=0.7, transform shape]
		    \umlclass[x=0, y=0]{MailingElement}{
		        -id: int
		    }{
		        +sendTo(): void
		    }
		
		    \umlclass[x=3, y=-4]{Email}{
		        -emailAddress: String \\
		        -firstName: String \\
		        -lastName: String \\
		    }{
		        +sendTo(): void \\
		    }
		
		    \umlclass[x=-4, y=-4]{MailingList}{
		        -emailAddress: String \\
		        -firstName: String \\
		        -lastName: String \\
		    }{
		        +sendTo(): void \\
		        +addMailingElement() \\
		        +deleteMailingElement() \\
		    }
		
		    \umlclass[x=8, y=-4]{MessengerContact}{
		        -contactId: String \\
		        -firstName: String \\
		        -lastName: String \\
		    }{
		        +sendTo(): void \\
		    }
		
		    \umlinherit{Email}{MailingElement}
		    \umlinherit{MessengerContact}{MailingElement}
		    \umlinherit{MailingList}{MailingElement}
		    \umlaggreg[geometry=|-, anchor1=north, anchor2=west]{MailingList}{Komponente}
		\end{tikzpicture}
	\end{center}
\end{frame}

\begin{frame}[fragile]{Komponente}
	\begin{exampleblock}{MailingElement.java}
 		\begin{lstlisting}
public abstract class MailingElement {
	private int id;
	
	public MailingElement(int id) {
		this.id = id;
	}
	
	public abstract void sendTo();
	
	public int getId() {
		return id;
	}
}    
        \end{lstlisting}
 	\end{exampleblock}
\end{frame}

\begin{frame}[fragile]{Blatt}
	\begin{exampleblock}{EmailAddress.java}
        \begin{lstlisting}
public class Email extends MailingElement {
	private final String firstName;
	private final String lastName;
	private final String emailAddress;
	
	public Email(int id, String firstName, 
            String lastName, String emailAddress) {
		super(id);
        // ...
	}
	
	@Override
	public void sendTo() {
		System.out.println("Email wird gesendet ");
		System.out.println("Empfaenger: " + firstName + " " 
            + lastName);
		System.out.println("Emailadresse: " + emailAddress);
	}            
        \end{lstlisting}
 	\end{exampleblock}
\end{frame}

\begin{frame}[fragile]{Noch ein Blatt}
	\begin{exampleblock}{MessengerContact.java}
        \begin{lstlisting}
public class MessengerContact extends MailingElement {
	private final String firstName;
	private final String lastName;
	private final String contactId;
	
	public MessengerContact(int id, String firstName, 
            String lastName, String contactId) {
		super(id);
		// ...
	}
	
	@Override
	public void sendTo() {
		System.out.println("Messengernachricht wird gesendet ");
		System.out.println("KontaktId: " + contactId);
	}            
        \end{lstlisting}
 	\end{exampleblock}
\end{frame}

\begin{frame}[fragile]{Kompositum}
	\begin{exampleblock}{MailingList.java}
        \begin{lstlisting}
public class MailingList extends MailingElement {
	private final List<MailingElement> mailingElements;
	private final String name;
	
	public MailingList(int id, String name) {
		super(id);
        // ...
	}
	
	@Override
	public void sendTo() {
		for (MailingElement mailingElement: mailingElements) {
			mailingElement.sendTo();
		}
	}
}            
        \end{lstlisting}
 	\end{exampleblock}
\end{frame}

\begin{frame}[fragile]{Kompositum}
	\begin{exampleblock}{MailingList.java}
        \begin{lstlisting}
public void addMailingElement(MailingElement mailingElement) {
    mailingElements.add(mailingElement);
}

public void deleteMailingElement(MailingElement mailingElement) {
    if (mailingElements.contains(mailingElement)) {
        mailingElements.remove(mailingElement);
    }
}
        \end{lstlisting}
 	\end{exampleblock}
\end{frame}

\begin{frame}[fragile]{Kompositum}
	\begin{exampleblock}{KompositumApplication.java}
        \begin{lstlisting}
private void deleteContact(List<MailingElement> mailingElements, 
        MailingElement contact) {
    for (MailingElement mailingElement: mailingElements) {
        if (mailingElement instanceof MailingList) {
            MailingList mailingList = 
                (MailingList)mailingElement;  
            mailingList.deleteMailingElement(contact);
        }
    }
}
        \end{lstlisting}
 	\end{exampleblock}
\end{frame}

\note{
    Wenn wir mit Polymorphie arbeiten, können wir bei diesem Muster auf ein Problem stoßen: Wir haben eine Liste mit der Basisklasse MailingElement und möchten die deleteMailingElement-Methode aufrufen. Das funktioniert allerdings nur, wenn wir eine Typumwandlung in die entsprechende Unterklasse vornehmen.
}

\begin{frame}{Allgemeines Klassendiagramm Alternative}
\begin{tikzpicture}[scale=0.7, transform shape]
    \umlclass[x=0, y=0]{Komponente}{
        -name: String
    }{
        +operation(): void \\
        \textit{+hinzufuegen(Komponente): void} \\
        \textit{+entfernen(Komponente): void} \\
        \textit{+getKind(int): Komponente} \\
    }

    \umlclass[x=4, y=-4]{Blatt}{
        -name: String
    }{
        +operation(): void\\
        +hinzufuegen(Komponente): void \\
        +entfernen(Komponente): void \\
        +getKind(int): Komponente \\
    }

    \umlclass[x=-4, y=-4]{Kompositum}{
        -name: String \\
        -elemente: List<Komponente> \\
    }{
        +operation(): void \\
        +hinzufuegen(Komponente): void \\
        +entfernen(Komponente): void \\
        +getKind(int): Komponente \\
    }

    \umlinherit{Blatt}{Komponente}
    \umlinherit{Kompositum}{Komponente}
    \umlaggreg[geometry=|-, anchor1=north, anchor2=west]{Kompositum}{Komponente}
\end{tikzpicture}
\end{frame}
\note{
	Diese Typumwandlungen könnte man vermeiden, wenn man die 
	Methoden zur Manipulation des Kompositums in die Oberklasse verschiebt und
	diese in den Blatt-Klassen leer überschreibt.
}

\begin{frame}{Klassendiagramm Organisation}
\begin{tikzpicture}[scale=0.7, transform shape]
    \umlclass[x=0, y=0]{OrganizationUnit}{
        -name: String
    }{
        +operation(): void \\
        \textit{+show(): void} \\
        \textit{+add(OrganizationUnit unit): void} \\
        \textit{+delete(OrganizationUnit unit): void} \\
    }

    \umlclass[x=4, y=-4]{Employee}{
        -name: String
    }{
        \textit{+show(): void} \\
        \textit{+add(unit: OrganizationUnit): void} \\
        \textit{+delete(unit: OrganizationUnit): void} \\
    }

    \umlclass[x=-4, y=-4]{Department}{
        -name: String \\
        -units: List<OrganizationUnit> \\
    }{
        \textit{+show(): void} \\
        \textit{+add(OrganizationUnit unit): void} \\
        \textit{+delete(OrganizationUnit unit): void} \\
    }

    \umlinherit{Employee}{OrganizationUnit}
    \umlinherit{Department}{OrganizationUnit}
    \umlaggreg[geometry=|-, anchor1=north, anchor2=west]{Department}{OrganizationUnit}
\end{tikzpicture}
\end{frame}
\note{
	Diese Typumwandlungen könnte man vermeiden, wenn man die 
	Methoden zur Manipulation des Kompositums in die Oberklasse verschiebt und
	diese in den Blatt-Klassen leer überschreibt.
}

\begin{frame}[fragile]{Komponente}
	\begin{exampleblock}{OrganizationUnit.java}
        \begin{lstlisting}
public abstract class OrganizationUnit {
    private String name;

    public OrganizationUnit(String name) {
        this.name = name;
    }

    public abstract void show();
    
    public abstract void add(OrganizationUnit unit);
    
    public abstract void delete(OrganizationUnit unit);
    
    public String getName() {
		return name;
	}
}
        \end{lstlisting}
 	\end{exampleblock}
\end{frame}

\begin{frame}[fragile]{Kompositum}
	\begin{exampleblock}{Department.java}
        \begin{lstlisting}
public class Department extends OrganizationUnit {
    private List<OrganizationUnit> units;

    @Override
    public void add(OrganizationUnit unit) {
        units.add(unit);
    }

    @Override
    public void delete(OrganizationUnit unit) {
    	units.remove(unit);
    }
}
        \end{lstlisting}
 	\end{exampleblock}
\end{frame}

\begin{frame}[fragile]{Blatt}
	\begin{exampleblock}{Employee.java}
        \begin{lstlisting}
public class Employee extends OrganizationUnit {
	@Override
	public void add(OrganizationUnit unit) {
	}

	@Override
	public void delete(OrganizationUnit unit) {
	}

    @Override
    public void show() {
        System.out.println("Mitarbeiter: " + getName());
    }
}
        \end{lstlisting}
 	\end{exampleblock}
\end{frame}

\begin{frame}{Klassendiagramm Dateisystem}
\begin{tikzpicture}[scale=0.7, transform shape]
    \umlclass[x=0, y=0]{FileSystemElement}{
        -name: String
    }{
        +operation(): void \\
        \textit{+display(level: int): void} \\
        \textit{+printFileNames(prefix: String): String} \\
        \textit{+isFile(): boolean} \\
    }

    \umlclass[x=4, y=-4]{FileElement}{
        -name: String
    }{
        +display(level: int): void \\
        +isFile(): boolean \\
        +printFileNames(prefix: String): String
    }

    \umlclass[x=-4, y=-4]{Directory}{
        -elements: List<FileSystemElement> \\
    }{
        +add(element: FileSystemElement): void \\
        +display(level: int): void \\
        +isFile(): boolean \\
        +printFileNames(prefix: String): String
    }

    \umlinherit{Employee}{OrganizationUnit}
    \umlinherit{Department}{OrganizationUnit}
    \umlaggreg[geometry=|-, anchor1=north, anchor2=west]{Department}{OrganizationUnit}
\end{tikzpicture}
\end{frame}
\note{
	Das nächste Beispiel liest ausgehend von einem bestimmten Verzeichnis alle Unterverzeichnisse und Dateien in eine Kompositum-Struktur zu leichteren Verwaltung der Dateien. 
}

\begin{frame}[fragile]{Komponente}
	\begin{exampleblock}{FileSystemElement.java}
        \begin{lstlisting}
public abstract class FileSystemElement {
    private String name;
    private String path;

    public FileSystemElement(String name, String path) {
        this.name = name;
        this.path = path;
    }

    public abstract void display(int level);
    public abstract boolean isFile();
    public abstract String printFileNames(String prefix);
    
    public String getFullName() {
        return path + File.separator + getName();
    }

    public String getName() {
        return name;
    }
}
        \end{lstlisting}
 	\end{exampleblock}
\end{frame}

\begin{frame}[fragile]{Blatt}
	\begin{exampleblock}{FileElement.java}
        \begin{lstlisting}
public class FileElement extends FileSystemElement {
    public FileElement(String name, String path) {
        super(name, path);
    }
    
    @Override
	public boolean isFile() {
        return true;
    }
    
    @Override
    public String printFileNames(String prefix) {
        return prefix + "/" + getName() + "\n";
    }
}
        \end{lstlisting}
 	\end{exampleblock}
\end{frame}

\begin{frame}[fragile]{Kompositum}
	\begin{exampleblock}{Directory.java}
        \begin{lstlisting}
public class Directory extends FileSystemElement {
    private final List<FileSystemElement> elements;

    public Directory(String name, String path) {
        super(name, path);
        elements = new ArrayList<>();
    }

    public void add(FileSystemElement element) {
        elements.add(element);
    }
    
    @Override
    public boolean isFile() {
        return false;
    }
        \end{lstlisting}
 	\end{exampleblock}
\end{frame}

\begin{frame}[fragile]{Kompositum}
	\begin{exampleblock}{Directory.java}
        \begin{lstlisting}
public void remove(String name) {
    Iterator<FileSystemElement> iterator = elements.iterator();
    while (iterator.hasNext()) {
        FileSystemElement element = iterator.next();
        if (element.getName().equals(name)) {
            iterator.remove();
            return;
        } else if (!element.isFile()) {
            ((Directory) element).remove(name);
        }
    }
}
        \end{lstlisting}
 	\end{exampleblock}
\end{frame}

\begin{frame}[fragile]{Kompositum}
	\begin{exampleblock}{Directory.java}
        \begin{lstlisting}
@Override
public String printFileNames(String prefix) {
    String currentPath = prefix + File.separator + getName();
    StringBuilder output = new StringBuilder();
    
    for (FileSystemElement element : elements) {
        output.append(element.printFileNames(currentPath)).
            append("\n");
    }
    
    return output.toString();
}
        \end{lstlisting}
 	\end{exampleblock}
\end{frame}